\section{Уравнение движения в механике сплошной среды}

Уравнение движения (баланса линейного импульса) в механике сплошной среды
(уравнение Ньютона для бесконечно малого объёма среды) имеет вид:
\begin{equation}
\rho \frac{\partial^2 u_i}{\partial t^2}
=
\frac{\partial \sigma_{ij}}{\partial x_j}.
\end{equation}

Здесь:
\begin{itemize}
    \item $\rho$ — плотность среды,
    \item $u_i(x_j,t)$ — $i$-я компонента вектора смещения,
    \item $\sigma_{ij}$ — тензор напряжений Коши,
    \item $x_j = (x_1,x_2,x_3)$ — пространственные координаты.
\end{itemize}

Левая часть уравнения описывает \textbf{инерцию} (плотность, умноженная на ускорение частиц среды),
правая часть — \textbf{дивергенцию тензора напряжений}, то есть внутренние силы.

\subsection{Смысл индексов}

\subsubsection{Индекс $i$}

\begin{itemize}
    \item свободный индекс,
    \item определяет направление уравнения,
    \item уравнение записывается для каждой компоненты:
\end{itemize}
\[
i = 1,2,3.
\]

\subsubsection{Индекс $j$}

\begin{itemize}
    \item повторяющийся индекс,
    \item по соглашению Эйнштейна происходит суммирование.
\end{itemize}

\[
\frac{\partial \sigma_{ij}}{\partial x_j}
=
\sum_{j=1}^{3}
\frac{\partial \sigma_{ij}}{\partial x_j}.
\]

Это выражение и есть \textbf{дивергенция тензора напряжений}.

\subsection{Запись по компонентам}

\subsubsection{Уравнение для $i=1$ (ось $x_1$)}

\begin{equation}
\rho \frac{\partial^2 u_1}{\partial t^2}
=
\frac{\partial \sigma_{11}}{\partial x_1}
+
\frac{\partial \sigma_{12}}{\partial x_2}
+
\frac{\partial \sigma_{13}}{\partial x_3}.
\end{equation}

\subsubsection{Уравнение для $i=2$ (ось $x_2$)}

\begin{equation}
\rho \frac{\partial^2 u_2}{\partial t^2}
=
\frac{\partial \sigma_{21}}{\partial x_1}
+
\frac{\partial \sigma_{22}}{\partial x_2}
+
\frac{\partial \sigma_{23}}{\partial x_3}.
\end{equation}

\subsubsection{Уравнение для $i=3$ (ось $x_3$)}

\begin{equation}
\rho \frac{\partial^2 u_3}{\partial t^2}
=
\frac{\partial \sigma_{31}}{\partial x_1}
+
\frac{\partial \sigma_{32}}{\partial x_2}
+
\frac{\partial \sigma_{33}}{\partial x_3}.
\end{equation}

Это система из трёх уравнений второго порядка по времени для компонент вектора смещений
\[
\mathbf{u} = (u_1,u_2,u_3).
\]

Каждое уравнение:
\begin{itemize}
    \item описывает движение вдоль соответствующей оси,
    \item учитывает все компоненты напряжений, действующие на элемент среды.
\end{itemize}

\subsection{Одномерный случай}

Если движение происходит только вдоль оси $x_1$:
\[
u_1 = u(x_1,t), \qquad u_2 = u_3 = 0,
\]
то система упрощается до
\begin{equation}
\rho \frac{\partial^2 u}{\partial t^2}
=
\frac{\partial \sigma_{11}}{\partial x_1},
\end{equation}
что соответствует обычному одномерному волновому уравнению.


\section{Термоупругое уравнение движения}

С учётом термоупругой связи уравнение движения записывается в виде:
\begin{equation}
\rho(x)\frac{\partial^2 u_i}{\partial t^2}
=
\frac{\partial}{\partial x_j}
\left[
\lambda(x)\delta_{ij}\frac{\partial u_k}{\partial x_k}
+
\mu(x)\left(
\frac{\partial u_i}{\partial x_j}
+
\frac{\partial u_j}{\partial x_i}
\right)
\right]
-
\gamma(x)\frac{\partial T}{\partial x_i}.
\end{equation}

\subsection{Индексы и суммирование}

\begin{itemize}
    \item $i = 1,2,3$ — индекс направления смещения (свободный),
    \item $j = 1,2,3$ — индекс пространственной производной (суммируемый),
    \item $k = 1,2,3$ — индекс суммирования (объёмная деформация).
\end{itemize}

\subsection{Неизвестные функции}

\begin{itemize}
    \item $u_i(x,t)$ — $i$-я компонента вектора смещения,
    \begin{itemize}
        \item физический смысл: перемещение частиц среды,
        \item размерность: $\mathrm{m}$,
    \end{itemize}
    \item $T(x,t)$ — температурное поле (отклонение от начальной температуры),
    \begin{itemize}
        \item размерность: $\mathrm{K}$.
    \end{itemize}
\end{itemize}

\subsection{Упругие параметры}

\begin{itemize}
    \item $\lambda(x)$ — первая константа Ламе,
    \begin{itemize}
        \item характеризует объёмное сжатие,
        \item связана с продольными (P) волнами,
        \item размерность: $\mathrm{Pa}$,
    \end{itemize}
    \item $\mu(x)$ — вторая константа Ламе (модуль сдвига),
    \begin{itemize}
        \item характеризует сопротивление сдвигу,
        \item определяет распространение поперечных (S) волн,
        \item размерность: $\mathrm{Pa}$.
    \end{itemize}
\end{itemize}

\subsection{Тензорные объекты}

Символ Кронкера:
\[
\delta_{ij}
=
\begin{cases}
1, & i=j, \\
0, & i \neq j.
\end{cases}
\]

\begin{itemize}
    \item выделяет диагональные компоненты,
    \item отражает изотропность среды.
\end{itemize}

\subsection{Термоупругая связь}

$\gamma(x)$ — коэффициент термоупругой связи, характеризующий влияние температурного градиента на механические напряжения.


\section*{Трёхмерная система уравнений термоупругости}

Рассмотрим декартову систему координат
\[
x_1 = x,\quad x_2 = y,\quad x_3 = z,
\]
и вектор перемещений
\[
\mathbf{u} = (u_1,u_2,u_3) = (u_x,u_y,u_z).
\]

Исходное уравнение движения термоупругой среды имеет вид:
\[
\rho \frac{\partial^2 u_i}{\partial t^2}
=
\frac{\partial}{\partial x_j}
\left[
\lambda \delta_{ij} \frac{\partial u_k}{\partial x_k}
+
\mu
\left(
\frac{\partial u_i}{\partial x_j}
+
\frac{\partial u_j}{\partial x_i}
\right)
\right]
-
\gamma \frac{\partial T}{\partial x_i}.
\]

\subsection*{Компонента $i=1$ (ось $x$)}

\begin{align}
\rho \frac{\partial^2 u_1}{\partial t^2}
&=
\frac{\partial}{\partial x}
\left[
\lambda
\left(
\frac{\partial u_1}{\partial x}
+
\frac{\partial u_2}{\partial y}
+
\frac{\partial u_3}{\partial z}
\right)
+
2\mu \frac{\partial u_1}{\partial x}
\right]
\nonumber\\
&\quad+
\frac{\partial}{\partial y}
\left[
\mu
\left(
\frac{\partial u_1}{\partial y}
+
\frac{\partial u_2}{\partial x}
\right)
\right]
+
\frac{\partial}{\partial z}
\left[
\mu
\left(
\frac{\partial u_1}{\partial z}
+
\frac{\partial u_3}{\partial x}
\right)
\right]
-
\gamma \frac{\partial T}{\partial x}.
\end{align}

\subsection*{Компонента $i=2$ (ось $y$)}

\begin{align}
\rho \frac{\partial^2 u_2}{\partial t^2}
&=
\frac{\partial}{\partial x}
\left[
\mu
\left(
\frac{\partial u_2}{\partial x}
+
\frac{\partial u_1}{\partial y}
\right)
\right]
\nonumber\\
&\quad+
\frac{\partial}{\partial y}
\left[
\lambda
\left(
\frac{\partial u_1}{\partial x}
+
\frac{\partial u_2}{\partial y}
+
\frac{\partial u_3}{\partial z}
\right)
+
2\mu \frac{\partial u_2}{\partial y}
\right]
\nonumber\\
&\quad+
\frac{\partial}{\partial z}
\left[
\mu
\left(
\frac{\partial u_2}{\partial z}
+
\frac{\partial u_3}{\partial y}
\right)
\right]
-
\gamma \frac{\partial T}{\partial y}.
\end{align}

\subsection*{Компонента $i=3$ (ось $z$)}

\begin{align}
\rho \frac{\partial^2 u_3}{\partial t^2}
&=
\frac{\partial}{\partial x}
\left[
\mu
\left(
\frac{\partial u_3}{\partial x}
+
\frac{\partial u_1}{\partial z}
\right)
\right]
+
\frac{\partial}{\partial y}
\left[
\mu
\left(
\frac{\partial u_3}{\partial y}
+
\frac{\partial u_2}{\partial z}
\right)
\right]
\nonumber\\
&\quad+
\frac{\partial}{\partial z}
\left[
\lambda
\left(
\frac{\partial u_1}{\partial x}
+
\frac{\partial u_2}{\partial y}
+
\frac{\partial u_3}{\partial z}
\right)
+
2\mu \frac{\partial u_3}{\partial z}
\right]
-
\gamma \frac{\partial T}{\partial z}.
\end{align}

